\songtitle{Three Realms, Four Rites}

\tag*{2022}\tag{medieval-epic}\tag{original-source}\tag{african-empire}
\begin{notes}
Part 2 of a five-part historical epic. Lyrics by Suno based on Carlos Thompson's historical summary of the \emph{African Empire} alternate-history universe.
\end{notes}
\begin{style}
symphonic medieval epic, Arabic maqam-inspired melodies, orchestral strings and low brass, oud and qanun, frame drums, darbuka, riqq, Byzantine-style choir, soaring mixed-language choral textures, solemn historical grandeur, cinematic dynamic arc, opening with oud and distant male choir, building through tense percussion and modal strings into a triumphant brass-and-choir climax, clear narrative male lead vocal, organic acoustic production with spacious stone-hall ambience, dramatic but dignified
\end{style}
\begin{lyrics}
\songpart{Intro}
From mountain passes, three banners rise,\\
Salt wind crossing the western skies.\\
A silver cross, a crescent flame,\\
Old roads remade, new crowns take name.

\songpart{Verse 1}
In the north, where the green hills lean,\\
Asturian bells ring over fields unseen.\\
Basque stone walls and Galician rain,\\
Guard the frontier, remember the slain.

Eastward, the old Roman tongue\\
Carries through vineyards, bright and young.\\
Occitan lords by the inland sea,\\
Build their halls where the roads run free.

\songpart{Pre-Chorus}
But southward sails a darker tide,\\
Black shields gathered at the desert side.\\
Through Maghreb gates the caravans turn,\\
While Moorish towers begin to burn.

\songpart{Chorus}
Three realms beneath one changing sun,\\
Three roads divided, three futures begun.\\
From Jerusalem’s stone to the Baltic snow,\\
Empires rise where the old winds blow.\\
O keepers of earth, O children of flame,\\
No border can silence the memory of names.

\songpart{Verse 2}
The horsemen thunder from the eastern plain,\\
Across Poland’s frost and Lithuania’s rain.\\
Their princes kneel beneath a different sign,\\
Cross and steppe in a single line.

Vilnius wakes with a conqueror’s crown,\\
From Arctic silence to the Black Sea town.\\
Urals behind them, forests ahead,\\
A Christian khanate where old worlds wed.

\songpart{Pre-Chorus}
In Avignon, the white-robed court\\
Makes a new Rome of the southern port.\\
Yet Rome remains, and Italy’s heart\\
Keeps its ancient fire, though realms divide apart.

\songpart{Chorus}
Four rites beneath one ancient throne,\\
Roman stone and Occitan tone.\\
Germanic bells, Anglican rain,\\
One distant shepherd through reform and change.\\
Jerusalem sings in a Maronite hall,\\
While Coptic lamps burn beyond the wall.

\songpart{Bridge}
At the Holy City’s eastern gate,\\
Crusaders arrive, but the years wait.\\
Byzantium guards the Anatolian height,\\
Avignon sends bread through winter night.\\
Maronite prayer, Ethiopian flame,\\
Mozarabic tongues preserve the name.\\
In Mali’s courts, the old rites stand,\\
Black kings blessing the desert land.

\songpart{Verse 3}
Baghdad falls beneath the mounted sky,\\
Turkish banners lift as the old courts die.\\
Mesopotamia changes its tongue and throne,\\
But Syria keeps the desert’s own.

From Arabia north to Africa’s shore,\\
The faithful kingdoms gather once more.\\
Yet Jerusalem’s walls and Anatolia’s stone\\
Hold Christian candles against the unknown.

\songpart{Chorus}
Three realms beneath one changing sun,\\
Four rites together, though the wars aren’t done.\\
From Maghreb sands to the Baltic snow,\\
Empires rise where the old winds blow.\\
O keepers of earth, O children of flame,\\
No border can silence the memory of names.

\annotation{Finale}\\
Then Swedish steel meets the steppe-born line,\\
And the sixteenth century redraws the sign.\\
The Mongol crown breaks, the provinces part,\\
But Vilnius remains in the continental heart.

Rome still burns, and Avignon prays,\\
Jerusalem keeps its guarded days.\\
Three realms, four rites, one world remade,\\
By every hand and every blade.

\songpart{Outro}
When the last bell fades over field and sea,\\
The past still chooses what the world will be.
\end{lyrics}

